% =====================================================================
\chapter{Variables continuas, desigualdades y ley de los grandes números}
\label{tema:6}
% =====================================================================

Muchas magnitudes, como tiempos, longitudes o pesos de una red neuronal, toman valores en un continuo; su distribución se describe con una densidad en lugar de una función de masa. Este tema introduce las variables continuas y tres distribuciones fundamentales (uniforme, exponencial y normal), y termina con las desigualdades de Markov y Chebyshev y la ley débil de los grandes números, que justifica que una media muestral aproxime una esperanza. Se siguen el capítulo~5 de Blitzstein y Hwang~\cite{bh} (§5.1, §5.2, §5.4 y §5.5) y su §10.1--10.2; las secciones §3.1 a §3.3 y §5.1 a §5.2 de Bertsekas y Tsitsiklis~\cite{bt} (clases~8 y~19 de MIT 6.041~\cite{mit6041f10}); los capítulos~5 y~8 de Ross~\cite{ross}; el capítulo~19 de~\cite{mcs} y el capítulo~8 de Grinstead y Snell~\cite{gs}.

\section{Densidad y función de distribución}

\begin{definicion}[Variable continua]
Una variable aleatoria $X$ es \emph{continua} si existe una función $f_X\colon\R\to[0,\infty)$, su \emph{densidad}, tal que para todos $a\le b$
\[
\PP(a\le X\le b)=\int_a^bf_X(x)\,dx.
\]
Necesariamente $\int_{-\infty}^{\infty}f_X(x)\,dx=1$. La \emph{función de distribución} es $F_X(x)=\PP(X\le x)=\int_{-\infty}^{x}f_X(t)\,dt$.
\end{definicion}

Consecuencias inmediatas: $\PP(X=x)=\int_x^xf_X=0$ para todo $x$, de modo que da igual escribir $<$ o $\le$; $F_X$ es continua y no decreciente, y en los puntos donde $f_X$ es continua, $F_X'(x)=f_X(x)$ por el teorema fundamental del cálculo. La densidad \emph{no} es una probabilidad y puede ser mayor que~1: lo que se interpreta es $\PP(x\le X\le x+\delta)\approx f_X(x)\,\delta$ para $\delta$ pequeño.

\begin{definicion}[Esperanza y varianza]
Para $X$ continua, $\E[X]=\int_{-\infty}^{\infty}x\,f_X(x)\,dx$ (si la integral converge absolutamente) y, para una función $g$, $\E[g(X)]=\int g(x)f_X(x)\,dx$. La varianza se define como en el caso discreto, $\Var(X)=\E[(X-\E X)^{2}]=\E[X^{2}]-\E[X]^{2}$, y la linealidad y las propiedades de la \cref{prop:t5-var} siguen siendo válidas.
\end{definicion}

\section{Distribuciones uniforme y exponencial}

\begin{definicion}[Uniforme]
$X\sim\Unif(a,b)$, con $a<b$, si $f_X(x)=\frac{1}{b-a}$ en $[a,b]$ y $0$ fuera. Entonces $F_X(x)=\frac{x-a}{b-a}$ para $x\in[a,b]$: la probabilidad de un subintervalo es proporcional a su longitud.
\end{definicion}

\begin{proposicion}\label{prop:t6-unif}
Si $X\sim\Unif(a,b)$, entonces $\E[X]=\frac{a+b}{2}$ y $\Var(X)=\frac{(b-a)^{2}}{12}$.
\end{proposicion}

\begin{proof}
$\E[X]=\frac1{b-a}\int_a^bx\,dx=\frac{b^{2}-a^{2}}{2(b-a)}=\frac{a+b}2$. Para la varianza basta el caso $\Unif(0,1)$, porque $X=a+(b-a)U$ con $U\sim\Unif(0,1)$ y $\Var(X)=(b-a)^{2}\Var(U)$. Para $U$, $\E[U^{2}]=\int_0^1u^{2}du=\frac13$ y $\Var(U)=\frac13-\frac14=\frac1{12}$.
\end{proof}

\begin{definicion}[Exponencial]
$X\sim\Expo(\lambda)$, con $\lambda>0$, si $f_X(x)=\lambda e^{-\lambda x}$ para $x\ge0$ y $0$ para $x<0$. Integrando, $F_X(x)=1-e^{-\lambda x}$ y $\PP(X>x)=e^{-\lambda x}$ para $x\ge0$. Se cumple $\E[X]=1/\lambda$ y $\Var(X)=1/\lambda^{2}$ (\cref{eje:6-expo}).
\end{definicion}

\begin{proposicion}[Falta de memoria]\label{prop:t6-memoria}
Si $X\sim\Expo(\lambda)$, entonces $\PP(X>s+t\mid X>s)=\PP(X>t)$ para todos $s,t\ge0$.
\end{proposicion}

\begin{proof}
Como $\{X>s+t\}\subseteq\{X>s\}$, $\PP(X>s+t\mid X>s)=\frac{e^{-\lambda(s+t)}}{e^{-\lambda s}}=e^{-\lambda t}$.
\end{proof}

La exponencial es la versión continua de la geométrica: modela tiempos de espera hasta un suceso que «no envejece», como la llegada de la siguiente petición a un servidor.

\section{Distribución normal y tipificación}

\begin{definicion}[Normal]
$Z\sim\Nor(0,1)$ (\emph{normal estándar}) si su densidad es $\varphi(z)=\frac{1}{\sqrt{2\pi}}e^{-z^{2}/2}$; su función de distribución se denota $\Phi$. Para $\mu\in\R$ y $\sigma>0$, $X\sim\Nor(\mu,\sigma^{2})$ si
\[
f_X(x)=\frac{1}{\sigma\sqrt{2\pi}}\exp\Bigl(-\frac{(x-\mu)^{2}}{2\sigma^{2}}\Bigr).
\]
\end{definicion}

Que $\varphi$ integra~1 equivale a la integral gaussiana $\int_{-\infty}^{\infty}e^{-z^{2}/2}dz=\sqrt{2\pi}$, que se demuestra en el módulo~5 con coordenadas polares. Por la simetría de $\varphi$, $\E[Z]=0$ y $\Phi(-z)=1-\Phi(z)$; integrando por partes se obtiene $\E[Z^{2}]=1$, luego $\Var(Z)=1$. $\Phi$ no tiene expresión con funciones elementales y se tabula (\cref{tab:t6-phi}).

\begin{teorema}[Tipificación]\label{teo:t6-tipificacion}
Si $Z\sim\Nor(0,1)$, entonces $X=\mu+\sigma Z\sim\Nor(\mu,\sigma^{2})$. Recíprocamente, si $X\sim\Nor(\mu,\sigma^{2})$, la variable tipificada $Z=(X-\mu)/\sigma$ es $\Nor(0,1)$ y
\[
\PP(X\le x)=\Phi\Bigl(\frac{x-\mu}{\sigma}\Bigr).
\]
En particular, $\E[X]=\mu$ y $\Var(X)=\sigma^{2}$.
\end{teorema}

\begin{proof}
Sea $X=\mu+\sigma Z$. Como $\sigma>0$, $F_X(x)=\PP(Z\le(x-\mu)/\sigma)=\Phi\bigl((x-\mu)/\sigma\bigr)$. Derivando con la regla de la cadena, $f_X(x)=\frac1\sigma\varphi\bigl(\frac{x-\mu}\sigma\bigr)$, que es la densidad $\Nor(\mu,\sigma^{2})$. El recíproco es el mismo cálculo leído al revés: si $X$ tiene esa densidad, $\PP\bigl((X-\mu)/\sigma\le z\bigr)=\PP(X\le\mu+\sigma z)=\int_{-\infty}^{\mu+\sigma z}\frac1\sigma\varphi\bigl(\frac{x-\mu}\sigma\bigr)dx=\int_{-\infty}^{z}\varphi(t)\,dt$ con el cambio $t=(x-\mu)/\sigma$. Por último, $\E[X]=\mu+\sigma\E[Z]=\mu$ y $\Var(X)=\sigma^{2}\Var(Z)=\sigma^{2}$.
\end{proof}

\begin{table}[h]
\centering
\begin{tabular}{@{}lcccccccc@{}}
\toprule
$z$ & $0{,}5$ & $1$ & $1{,}5$ & $1{,}645$ & $1{,}96$ & $2$ & $2{,}5$ & $3$ \\
\midrule
$\Phi(z)$ & $0{,}6915$ & $0{,}8413$ & $0{,}9332$ & $0{,}9500$ & $0{,}9750$ & $0{,}9772$ & $0{,}9938$ & $0{,}9987$ \\
\bottomrule
\end{tabular}
\caption{Valores de la función de distribución normal estándar. Para $z<0$, $\Phi(z)=1-\Phi(-z)$.}
\label{tab:t6-phi}
\end{table}

\begin{ejemplo}[Regla 68--95--99,7]\label{ej:t6-regla}
Si $X\sim\Nor(\mu,\sigma^{2})$, $\PP(\abs{X-\mu}\le k\sigma)=\PP(\abs Z\le k)=2\Phi(k)-1$, que vale $0{,}6827$, $0{,}9545$ y $0{,}9973$ para $k=1,2,3$. Por ejemplo, si las estaturas siguen una $\Nor(170,\,10^{2})$ (en cm), $\PP(\text{estatura}>185)=1-\Phi(1{,}5)=0{,}0668$.
\end{ejemplo}

\section{Desigualdades de Markov y Chebyshev}

\begin{teorema}[Desigualdad de Markov]\label{teo:t6-markov}
Si $X\ge0$ y $a>0$, entonces $\PP(X\ge a)\le\dfrac{\E[X]}{a}$.
\end{teorema}

\begin{proof}
Se cumple la desigualdad entre variables $a\,I_{\{X\ge a\}}\le X$: si $X\ge a$, el lado izquierdo vale $a\le X$; si no, vale $0\le X$. Tomando esperanzas, que conservan las desigualdades por linealidad (si $Y\ge0$, $\E[Y]\ge0$), $a\,\PP(X\ge a)\le\E[X]$.
\end{proof}

\begin{teorema}[Desigualdad de Chebyshev]\label{teo:t6-chebyshev}
Si $X$ tiene media $\mu$ y varianza $\sigma^{2}$ finitas, para todo $c>0$,
\[
\PP(\abs{X-\mu}\ge c)\le\frac{\sigma^{2}}{c^{2}}.
\]
\end{teorema}

\begin{proof}
Se aplica Markov a la variable no negativa $(X-\mu)^{2}$ con $a=c^{2}$, observando que $\{\abs{X-\mu}\ge c\}=\{(X-\mu)^{2}\ge c^{2}\}$: $\PP(\abs{X-\mu}\ge c)\le\E[(X-\mu)^{2}]/c^{2}=\sigma^{2}/c^{2}$.
\end{proof}

Con $c=k\sigma$, Chebyshev dice que $\PP(\abs{X-\mu}\ge k\sigma)\le1/k^{2}$ para \emph{cualquier} distribución: a dos desviaciones típicas hay como mucho un $25\,\%$ de la probabilidad. Para la normal el valor exacto es un $4{,}6\,\%$; la cota es peor porque no usa nada de la forma de la distribución.

\section{Ley débil de los grandes números}

\begin{proposicion}[Media muestral]\label{prop:t6-media}
Sean $X_1,\dots,X_n$ independientes con la misma media $\mu$ y la misma varianza $\sigma^{2}$, y sea $\Xbar_n=\frac1n\sum_{i=1}^{n}X_i$. Entonces $\E[\Xbar_n]=\mu$ y $\Var(\Xbar_n)=\sigma^{2}/n$.
\end{proposicion}

\begin{proof}
Por linealidad, $\E[\Xbar_n]=\frac1n\cdot n\mu=\mu$. Por la \cref{prop:t5-producto} (válida también para variables continuas independientes) y la \cref{prop:t5-var}, $\Var(\Xbar_n)=\frac1{n^{2}}\sum_i\Var(X_i)=\frac{n\sigma^{2}}{n^{2}}$.
\end{proof}

\begin{teorema}[Ley débil de los grandes números]\label{teo:t6-lgn}
En las condiciones de la \cref{prop:t6-media}, para todo $\varepsilon>0$,
\[
\PP\bigl(\abs{\Xbar_n-\mu}\ge\varepsilon\bigr)\le\frac{\sigma^{2}}{n\varepsilon^{2}}\xrightarrow[n\to\infty]{}0.
\]
Se dice que $\Xbar_n$ \emph{converge en probabilidad} a $\mu$.
\end{teorema}

\begin{proof}
Se aplica la desigualdad de Chebyshev a $\Xbar_n$, cuya media es $\mu$ y cuya varianza es $\sigma^{2}/n$.
\end{proof}

\begin{corolario}[Frecuencias relativas]\label{cor:t6-frecuencias}
Si $A$ es un suceso de probabilidad $p$ y $f_n$ es la frecuencia relativa con que ocurre en $n$ repeticiones independientes, entonces $\PP(\abs{f_n-p}\ge\varepsilon)\le\dfrac{p(1-p)}{n\varepsilon^{2}}\le\dfrac{1}{4n\varepsilon^{2}}$.
\end{corolario}

\begin{proof}
$f_n$ es la media de $n$ indicadoras independientes $\Ber(p)$, de varianza $p(1-p)$, y $p(1-p)=\frac14-\bigl(p-\frac12\bigr)^{2}\le\frac14$.
\end{proof}

El corolario reconcilia la definición axiomática con la interpretación frecuentista: la probabilidad de un suceso es el valor al que tiende, en probabilidad, su frecuencia relativa. Con $\varepsilon=0{,}01$ y $n=10^{5}$, la cota es $1/(4\cdot10^{5}\cdot10^{-4})=0{,}025$.

\begin{observacion}[Conexión con el curso]
La normal aparece en la inicialización de los pesos (módulo~5) y en la activación GELU, $x\,\Phi(x)$ (módulo~11). La ley de los grandes números justifica que la entropía cruzada media sobre un corpus estime la pérdida esperada (módulos~1 y~5) y que el gradiente de un minilote estime el gradiente completo (módulo~7); la \cref{prop:t6-media} explica que el error típico de estas estimaciones decrezca como $1/\sqrt n$. El módulo~5 completa este tema con el teorema central del límite.
\end{observacion}

\begin{observacion}[Errores frecuentes]
(1) Interpretar $f_X(x)$ como una probabilidad o exigir que sea $\le1$. (2) Olvidar que $\PP(X=x)=0$ para variables continuas, o escribir que un suceso de probabilidad~0 es imposible. (3) Tipificar con la varianza en lugar de la desviación típica: $Z=(X-\mu)/\sigma$, no $(X-\mu)/\sigma^{2}$. (4) Aplicar Markov a variables que pueden ser negativas. (5) Leer la ley de los grandes números como una «compensación»: tras muchas caras seguidas, la probabilidad de cruz sigue siendo $1/2$; la frecuencia se acerca a $1/2$ porque las desviaciones pasadas quedan diluidas, no corregidas.
\end{observacion}

\begin{lecturas}
\item Blitzstein y Hwang~\cite{bh}, cap.~5, §5.1 «Probability density functions», §5.2 «Uniform», §5.4 «Normal» y §5.5 «Exponential»; cap.~10, §10.1 «Inequalities» y §10.2 «Law of large numbers».
\item Bertsekas y Tsitsiklis~\cite{bt}, §3.1 a §3.3 y §5.1 a §5.3 (clases~8 y~19 de MIT 6.041~\cite{mit6041f10}).
\item Ross~\cite{ross}, cap.~5, §5.1 a §5.5, y cap.~8, §8.2 «Chebyshev's Inequality and the Weak Law of Large Numbers».
\item Lehman, Leighton y Meyer~\cite{mcs}, cap.~19 «Deviation from the Mean», §19.1 a §19.3.
\item Grinstead y Snell~\cite{gs}, §2.2 «Continuous Density Functions», §5.2 «Important Densities» y cap.~8 «Law of Large Numbers».
\end{lecturas}

\section*{Ejercicios del tema 6}
\addcontentsline{toc}{section}{Ejercicios del tema 6}

\begin{ejercicio}{1}{eje:6-conceptos}
Justifique si son verdaderas: (a) una densidad no puede tomar valores mayores que~1; (b) si $X$ es continua, $\PP(X\le 2)=\PP(X<2)$; (c) si $\PP(A)=0$, el suceso $A$ es imposible; (d) toda función de distribución es no decreciente.
\end{ejercicio}

\begin{ejercicio}{2}{eje:6-densidad}
Sea $f(x)=cx^{2}$ para $0\le x\le2$ y $f(x)=0$ en otro caso. (a) Halle $c$ para que $f$ sea una densidad. (b) Halle la función de distribución. (c) Calcule $\PP(X>1)$, $\E[X]$ y $\Var(X)$.
\end{ejercicio}

\begin{ejercicio}{2}{eje:6-pasas}
En una caja de cereales de 30\,cm de altura, la densidad de pasas a la altura $h$ (pasas por centímetro) es $40-h$. (a) ¿Cuántas pasas hay en la caja? (b) Sea $H$ la altura de una pasa elegida al azar. Halle su densidad y su función de distribución. (c) ¿Qué probabilidad hay de que esté en el tercio inferior de la caja?
\fuente{adaptado de MIT 18.05~\cite{mit1805}, \emph{Problem Set}~3, problema~4.}
\end{ejercicio}

\begin{ejercicio}{2}{eje:6-exponencial}
Sea $Y\sim\Expo(\lambda)$. (a) Obtenga la función de distribución integrando la densidad. (b) Calcule $\PP(Y\le1/\lambda)$, es decir, la probabilidad de que $Y$ no supere su media. (c) Halle la mediana de $Y$.
\fuente{adaptado de MIT 18.05~\cite{mit1805}, \emph{Problem Set}~3, problema~5(c).}
\end{ejercicio}

\begin{ejercicio}{2}{eje:6-autobus}
En la ciudad A pasa un autobús exactamente cada 10 minutos, y un viajero llega en un instante uniforme, de modo que su espera es $\Unif(0,10)$. En la ciudad B la espera es exponencial de media 10 minutos. (a) Calcule la espera media en cada ciudad. (b) En cada ciudad, si el viajero ya lleva 6 minutos esperando, ¿qué probabilidad hay de que espere al menos 3 minutos más?
\fuente{adaptado de Harvard Stat 110~\cite{stat110}, \emph{Homework}~6, problema~1, apartados (a) a~(c).}
\end{ejercicio}

\begin{ejercicio}{2}{eje:6-tipificar}
Sea $X\sim\Nor(2,3^{2})$. Usando la \cref{tab:t6-phi}, calcule $\PP(X\le2)$, $\PP(X-\mu>1{,}5\sigma)$ y $\PP(\abs{X-\mu}<1{,}5\sigma)$.
\fuente{adaptado de MIT 18.05~\cite{mit1805}, \emph{Problem Set}~3, problema~5(b).}
\end{ejercicio}

\begin{ejercicio}{2}{eje:6-gestacion}
La duración de un embarazo, en días, se modela como $\Nor(280,\,8{,}5^{2})$. (a) ¿Qué probabilidad hay de que el parto ocurra al menos una semana antes de la fecha prevista (día 280), es decir, como muy tarde el día 273? (b) ¿Qué día $d$ cumple que el parto ocurre después de $d$ con probabilidad $0{,}95$?
\fuente{adaptado de MIT 18.05~\cite{mit1805}, \emph{Problem Set}~3, problema~6 (con otras preguntas).}
\end{ejercicio}

\begin{ejercicio}{2}{eje:6-palo}
Se rompe un palo de longitud~1 por un punto $U\sim\Unif(0,1)$. Sea $L$ la longitud del trozo más largo. Halle la función de distribución de $L$, identifique su distribución y calcule $\E[L]$.
\fuente{adaptado de Harvard Stat 110~\cite{stat110}, \emph{Strategic Practice}~5, sección «Continuous Distributions», problema~4.}
\end{ejercicio}

\begin{ejercicio}{2}{eje:6-markov}
Sea $X\sim\Expo(1)$. (a) Acote $\PP(X\ge3)$ con la desigualdad de Markov y compare con el valor exacto. (b) Acote $\PP(\abs{X-1}\ge2)$ con la de Chebyshev y compare con el valor exacto.
\end{ejercicio}

\begin{ejercicio}{2}{eje:6-n-cheb}
Sean $X_1,\dots,X_n$ independientes con media $\mu$ y varianza $\sigma^{2}$. Use la desigualdad de Chebyshev para hallar un $n$ que garantice que la media muestral dista de $\mu$ menos de $2\sigma$ con probabilidad al menos $0{,}99$.
\fuente{adaptado de Harvard Stat 110~\cite{stat110}, \emph{Strategic Practice}~10, sección «Inequalities», problema~3.}
\end{ejercicio}

\begin{ejercicio}{2}{eje:6-unif-lgn}
Sean $U_1,U_2,\dots$ independientes $\Unif(-1,1)$ y $\overline U_n$ su media muestral. (a) Demuestre con la ley débil de los grandes números que $\overline U_n$ converge en probabilidad a~0. (b) Acote $\PP(\abs{\overline U_{100}}\ge0{,}1)$ con Chebyshev.
\fuente{adaptado de MIT 6.041~\cite{mit6041f10}, \emph{Problem Set}~9, problema~1(a).}
\end{ejercicio}

\begin{ejercicio}{3}{eje:6-encuesta}
En una encuesta a $n$ personas elegidas independientemente, cada una apoya una propuesta con probabilidad $p$ desconocida, y $\hat p$ es la proporción de apoyos en la muestra. (a) Demuestre que $\PP(\abs{\hat p-p}>c)\le\frac{1}{4nc^{2}}$ para todo $c>0$. (b) ¿Qué tamaño de muestra garantiza, sea cual sea $p$, un error menor que $0{,}05$ con probabilidad al menos $0{,}95$?
\fuente{adaptado de Harvard Stat 110~\cite{stat110}, \emph{Penultimate Homework}, problema~6.}
\end{ejercicio}

\begin{ejercicio}{3}{eje:6-cheb-exacta}
Dados $\mu\in\R$, $\sigma>0$ y $c\ge\sigma$, construya una variable aleatoria discreta con tres valores, media $\mu$ y desviación típica $\sigma$, para la que la desigualdad de Chebyshev $\PP(\abs{X-\mu}\ge c)\le\sigma^{2}/c^{2}$ sea una igualdad. Concluya que la cota no puede mejorarse sin más hipótesis.
\fuente{adaptado de MIT 6.041~\cite{mit6041f10}, \emph{Problem Set}~9, problema~2.}
\end{ejercicio}

\begin{ejercicio}{3}{eje:6-cantelli}
Sea $X$ con media $\mu$ y varianza $\sigma^{2}$, y $a>0$. Demuestre la desigualdad de Chebyshev unilateral (de Cantelli): $\PP(X-\mu\ge a)\le\dfrac{\sigma^{2}}{\sigma^{2}+a^{2}}$. \emph{Indicación:} para $t\ge0$, $\PP(X-\mu\ge a)\le\PP\bigl((X-\mu+t)^{2}\ge(a+t)^{2}\bigr)$; aplique Markov y elija $t$ de forma óptima.
\fuente{adaptado de MIT 6.041~\cite{mit6041f10}, \emph{Problem Set}~9, problema~G1.}
\end{ejercicio}

\begin{ejercicio}{3}{eje:6-expo}
Sea $X\sim\Expo(\lambda)$. Demuestre, integrando por partes, que $\E[X]=1/\lambda$ y $\E[X^{2}]=2/\lambda^{2}$, y deduzca $\Var(X)$. Demuestre además que $\lambda X\sim\Expo(1)$.
\end{ejercicio}

\begin{ejercicio}{3}{eje:6-afin-normal}
Sea $X\sim\Nor(\mu,\sigma^{2})$ y $a,b\in\R$ con $a\neq0$. Demuestre que $aX+b\sim\Nor(a\mu+b,\,a^{2}\sigma^{2})$, tratando por separado los casos $a>0$ y $a<0$. Deduzca que si $Z\sim\Nor(0,1)$, también $-Z\sim\Nor(0,1)$.
\end{ejercicio}
